% DRAFT for IGE/paperA_joint_loss/sec5_household.tex, subsection "The safety net" (WP6, 2026-10-20).
% Replaces the current one-paragraph subsection. Not yet in the paper: the lead merges it (do not \input this file).
% Numbers: IGE/78.SafetyNetJointLoss.py -> IGE/output/safety_net_joint_loss.{csv,md,png}.
% The existing PSID event-study numbers (food stamps +4.2 / +3.3 points) are from IGE/19 via IGE/60 and are kept.

\subsection{The safety net}

Food stamp receipt rises by 4.2 points for Black and 3.3 points for White families after the father's displacement,
while AFDC/TANF receipt barely moves. The joint job loss determines which programs can respond, and the monthly SIPP
shows how. When the husband loses his job alone, unemployment insurance reaches him at once: his receipt rises by 23
points in the first three months and stays 13 points higher half a year later, and the couple almost never goes
without both earnings and benefits (1.0\% of couple-months in the first half year). When the wife loses her job too,
his benefits rise by about as much, hers rise by 11 to 18 points, and yet 9.3\% of the couple-months in the first half
year have no earner and no unemployment insurance, 5.7 points more than after his loss alone.\footnote{SIPP panels
2004, 2008 and 2014 (reference months 2003--2016): dual-earner couples with the wife aged 25--59, both employed in
the month before the husband's involuntary separation (layoff, closure, discharge or end of a temporary job) at
$s = 0$; a joint loss is an involuntary separation of the wife in $s = -3, \dots, +3$ (206 events, against 2{,}644
losses of his alone). Clean controls: 29{,}899 couples with a pseudo-event month in which both work and neither
spouse separates involuntarily nearby. Regressions on event month $\times$ group with couple and panel $\times$
rotation $\times$ calendar-month fixed effects, reference $s = -3, \dots, -1$, standard errors clustered by household
(32{,}109 clusters, $N = 516{,}624$). His UI receipt: $+22.6^{***}$ points (s.e.\ 0.7) in $s = 0$--$2$ and
$+13.3^{***}$ (0.7) in $s = 6$--$11$ after his loss alone; $+18.2^{***}$ (2.5) and $+9.7^{***}$ (2.5) after a joint
loss. Her UI receipt after a joint loss: $+10.9^{***}$ (2.1) in $s = 0$--$2$ and $+18.0^{***}$ (2.7) in $s = 3$--$5$.
No earner and no UI: $+6.2^{***}$ (1.4) and $+6.4^{***}$ (1.7) after a joint loss, $+0.5^{***}$ (0.1) and
$+0.7^{***}$ (0.2) after his loss alone; joint minus alone $+5.7^{***}$ (1.7; $p<0.001$) in $s = 3$--$5$ and
$+2.6^{*}$ (1.3; $p=0.055$) in $s = 6$--$11$. Receipt is from the core-wave program items (2004/08: state UI ER05 /
T05AMT, household food stamps THFDSTP, own receipt and coverage ER27, RCUTYP27; 2014: EUC1--3MNYN, RFSYN, RFSCOV).
The employment pre-trend ($-5.9^{***}$ in $s = -6$ to $-4$) follows from requiring both spouses to work at $s = -1$
(\texttt{IGE/78.SafetyNetJointLoss.py}).} In the couple-months of a joint loss in which neither spouse works, which
are a quarter of the first half year, 39\% receive unemployment insurance for neither spouse and 23\% receive neither
unemployment insurance nor food stamps.\footnote{Couple-months $s = 1$--$6$ after a joint loss with both spouses
not employed: 24.0\% of the joint-loss couple-months (1.5\% after his loss alone); among them UI for neither spouse
38.7\%, food stamps 33.8\%, neither program 23.5\% (316 couple-months). After his loss alone, among couple-months with
him not employed: UI for neither 51.2\%, food stamps 9.7\% (6{,}458 couple-months). Weighted by the wife's weight in
$s = -1$ (\texttt{IGE/78.SafetyNetJointLoss.py}).}

Food stamps respond less than the joint loss would suggest, because the couples who lose both jobs are already the
couples on the program. In the three months before his loss, 19\% of joint-loss couples receive food stamps, against
5.9\% of couples in which he loses his job alone and 3.2\% of the controls; receipt then rises by 1.0 to 1.5 points after
his loss alone and by 1 to 3 points, imprecisely, after a joint loss. AFDC/TANF does not move in either
group.\footnote{Food stamp receipt relative to $s = -3, \dots, -1$: after his loss alone $+1.05^{***}$ (0.26) in
$s = 0$--$2$, $+1.51^{***}$ (0.36) in $s = 3$--$5$, $+1.40^{***}$ (0.43) in $s = 6$--$11$; the amount rises by
\$3.9--5.2 a month (2000 dollars, zeros included). After a joint loss $+1.14$ (1.25; $p=0.361$), $+3.28$ (2.23;
$p=0.141$), $+3.26$ (2.61; $p=0.212$); joint minus alone $+1.77$ (2.26; $p=0.434$) in $s = 3$--$5$. TANF receipt:
$+0.12$ (0.08; $p=0.138$) and $+0.16$ (0.58; $p=0.779$) in $s = 0$--$2$. In the PSID couple-years (25's intervals,
1984--2021), food stamps are received in the year before by 15.6\% of joint-loss couples, 10.5\% after his
displacement alone and 3.4\% when neither spouse is displaced; the change over the interval relative to couples with
no displacement is $+1.88^{***}$ points (0.57) after his displacement alone and $+5.24^{**}$ (2.18; $p=0.016$) after a
joint one, and AFDC/TANF does not change ($+0.17$, $p=0.571$; $+0.11$, $p=0.932$); calendar-year fixed effects, 1968
family clusters, $N = 60{,}061$ (\texttt{IGE/78.SafetyNetJointLoss.py}).}

An earnings subsidy such as the Earned Income Tax Credit pays nothing in a month in which neither spouse works, and
these are the months the joint loss adds. Over the calendar year the schedule acts the other way. Before the loss,
most dual-earner couples earn too much for the credit, and the fall in annual earnings moves some of them into its
range. Applying the married-joint schedule of the era to the PSID couples' earnings, 24\% of joint-loss couples are in
the credit range in the year before and 34\% at the end of the interval. 8.3\% end in the phase-in region or with no
earnings, where the credit falls with every lost dollar, against 3.8\% after his displacement alone. The schedule
credit rises by \$196 a year relative to couples without a displacement, against \$66 after his displacement alone.
The subsidy pays least to the couples whose joint loss lasts long enough to empty a year of earnings. For the others
it is a partial, delayed insurance against the loss of annual income.\footnote{Schedule of
\texttt{IGE/61.PolicySchedules.py} (era means, 2000 dollars, married filing jointly, children capped at three),
applied to the sum of both spouses' annual labor incomes; no AGI or investment-income test, no take-up, so this
places the couple on the schedule and does not measure receipt. Year before = the interview opening the interval,
end = the interview closing it (one year to 1996, two after). Shares at the end of the interval for joint-loss
couples: no earnings 0.9\%, phase-in 7.4\%, plateau 4.6\%, phase-out 21.2\%, above the range 66.0\% (326 couple-years).
LPMs on displacement groups relative to couples without a displacement, calendar-year fixed effects, 1968 family
clusters, $N = 60{,}061$: change in the schedule credit $+\$195.7^{***}$ (48.8) joint, $+\$65.8^{***}$ (11.6) his
alone, difference $+\$129.9^{***}$ (49.1; $p=0.008$); entering the phase-in region or no earnings $+4.86^{***}$
points (1.36) joint, $+1.69^{***}$ (0.36) his alone, difference $+3.17^{**}$ (1.40; $p=0.023$). Black couples: credit
$+\$287^{**}$ (115) joint, entering phase-in or no earnings $+8.33^{***}$ (2.92); White couples: $+\$153^{***}$ (45)
and $+3.10^{**}$ (1.42). Annual intervals only (1984--1996): credit $+\$141^{***}$ (39) joint, $+\$55^{***}$ (10) his
alone (\texttt{IGE/78.SafetyNetJointLoss.py}).} Unemployment insurance covers each spouse only if he or she
qualifies, so a joint loss can leave one of the two uncovered: in nearly four in ten joint-loss months without an earner,
neither spouse draws it.

% ---------------------------------------------------------------------------------------------------------------
% sec7 'Counterfactuals' (iii), motivation sentence (draft):
% "The ranking rests on the timing of the two instruments: in the first half year after a joint loss, 9.3% of
%  couple-months have no earner and no unemployment insurance (1.0% after his loss alone), and in those months the
%  earnings subsidy pays nothing, while over the year the schedule credit of joint-loss couples rises by $196
%  (Section 5)."
% NOTE for the PI: at annual frequency the claim 'pays least when both spouses have lost their jobs' holds only for
%  the 8% of joint-loss couples who end in the phase-in region or with no earnings; on average the annual credit
%  RISES after a joint loss (+$196***), because 76% of these couples are above the range before. In the monthly
%  model, the counterfactual should apply the credit to annual earnings if the comparison is to match the statute.
% ---------------------------------------------------------------------------------------------------------------
% Validation rows for IGE/paperA_model/fit_table.py (untargeted; do not add to the objective):
%  sipp_snap_joint_s16   SNAP receipt, joint-loss couple-months s = 1..6 (SIPP, weighted)      21.6%  (N 1,200)
%  sipp_snap_alone_s16   SNAP receipt, his-loss-alone couple-months s = 1..6                   7.45%  (N 15,352)
%  sipp_noui_joint_both  UI for neither spouse, joint-loss months with neither employed       38.7%  (N 316)
%  sipp_noinc_joint_s16  no earner and no UI, joint-loss couple-months s = 1..6                9.28%  (N 1,200)
%  sipp_noinc_alone_s16  no earner and no UI, his-loss-alone couple-months s = 1..6           1.01%  (N 15,352)
%  psid_eitc_lowreg_joint  share in phase-in or no earnings at the end of the interval, joint    8.3%  (N 326)
